\section{Conclusions}
\label{sec:conclusions}

\paragraph{Findings.}
Conformal prediction turns frozen monocular metric depth models into per-pixel interval estimators with finite-sample guarantees, provided the guarantee is stated for the right unit: the image or the scene, not the pixel.
A cross-image guarantee is valid at modest cost, but it is an average that hides images with large global scale errors, and it does not survive a change of dataset.
A few sparse metric anchors remove most of the scale error (intervals 1.1--3.7$\times$ narrower).
Calibrating inside each image on its own anchors gives a per-image guarantee only for uniformly random anchors; under the layouts real sensors produce, it fails silently, with tighter intervals that miss most of the image.
Giving the sensor budget two roles fixes this.
Structured anchors build the prediction, and a few uniformly random probes, hidden from the predictor, calibrate it within each image.
The resulting layer needs no calibration images for validity (held-out images only tune its width), turns the probe count into a design parameter with a known per-image coverage law, and runs on any probe-independent centre.
A depth-completion network is the better centre under uniform anchors, and the frozen model with an anchor-built centre under clustered anchors (ten confirmatory splits, one network, one dataset).
Laser probes certified the depth even with iPad anchors; iPad probes certified only the iPad reading.

\paragraph{Limitations.}
\begin{itemize}[leftmargin=*,itemsep=1pt,topsep=2pt]
\item \emph{Expectation, not every draw:} with 20 probes about \NumBetaShortOneTwenty\% of images fall below 0.85 coverage in a single draw.
\item \emph{No spatial guarantee from probes alone:} the per-distance guarantee of \mDualBoth needs calibration scenes and the same anchor rule at deployment.
\item \emph{What is certified:} the probes' quantity.
The envelope that converts reading coverage into depth coverage holds only on average over scenes and, in our runs, costs more than \mDistBoth.
ARKit depth is fused and dense; raw sparse returns and dropouts on glass or sky are untested.
\item \emph{Selection of $\gamma$ and the design of \mDist} used NYU and KITTI, so only ARKitScenes checks efficiency on held-out data, and UniDepthV2 saw ARKitScenes in training.
\item \emph{Validity needs no images; efficiency does:} the \mDist variant, floor $\tau$ and distance function $g$ are chosen on 200 head-train images with dense ground truth under the same anchor layout. Without such images \cref{thm:dual} still holds, but the intervals may be wider.
\item \emph{Scope:} one $\alpha$ in the main method, two budgets, one independent sensor and uniform probes by design.
With a fixed-pattern scanner the guarantee covers only pixels with a return.
There is no downstream evaluation beyond the safety-query study.
\item \emph{Efficiency under clustered anchors:} there \mDual's interval score is up to \NumDualISLoseMax\% above the hull baseline's; its advantage is the per-image guarantee, not the score.
\end{itemize}

\paragraph{Future work.}
\begin{itemize}[leftmargin=*,itemsep=1pt,topsep=2pt]
\item A tighter sensor-error envelope that does not pay for both error sources separately.
\item Test raw, sparse sensor returns instead of fused depth, and video, where the probes could be shared across frames.
\item Study active placement of the anchors (the probes must stay random).
\item Evaluate the intervals inside a downstream AR or planning task, calibrating the guarantee to the query it protects, as the safety-query study suggests.
\end{itemize}
